\chapter{Metodología y Diseño del Sistema}
\label{cap:MetodologiaDisenoSistema}

\section{Generación del Dataset}
El sistema genera un número (N) de muestras de audio con valores continuos aleatorios,
en base a unos rangos de valores específicos para cada parámetro. Para cada uno de ellos 
se utiliza una distribución uniforme, a excepción de la frecuencia portadora (\textit{Carrier}) que 
se realiza en escala logarítmica. De esta forma se simula la percepción humana del sonido. 

En primer lugar, se sintetizan los N archivos \textit{.wav} mediante la librería \textbf{NumPy} y la formulación matemática de la FM (Fase 1). 
Se define una frecuencia de muestreo (\textit{sample rate}) de 44100 Hz
y se incrementa el tiempo de duración de 1 a 2 segundos, con el fin de asegurar un margen temporal
necesario para analizar y percibir las variaciones en el sonido provocadas por las envolventes. 
Estas envolventes se generan sobre un vector de tiempo unificado. 

En la Fase 2, los audios generados \textit{.wav} se convierten a tensores nativos de \textit{PyTorch}. Se cargan mediante \textit{torchaudio} y
se les aplica un suavizado de amplitud en los extremos (\textit{fade-in} y \textit{fade-out}) con el objetivo de 
evitar chasquidos o ruidos indeseados. Tras una normalización de pico, se aplican las dos transformaciones principales:
de audio a espectrograma y de amplitud lineal a logarítmica (dB). Para la representación acustica, el programa permite elegir
entre dos métodos, con el fin de llevar a cabo diferentes pruebas de experimentación. Por un lado, se encuentran 
los espectrogramas en formato lineal, mediante \textbf{STFT} con 513  \textit{bins}\footnote{El término \textit{bin} hace referencia los intervalos de la Tranformada de Fourier (FFT). 
El tamaño de ventana utilizado para la creación de los espectrogramas es de 1024 ($N_{fft} = 1024$), lo que genera
$1024 / 2 + 1 = 513$ intervalos o \textit{bins} de frecuencias}de frecuencia. Y por el otro lado, 
se da lugar a elegir un formato logarítmico-perceptual (\textbf{Mel} con 128 bandas), el cual tiene un mayor potencial para
la representación del sonido de manera similar a la percibida por el humano.

\section{Preprocesamiento de Datos: normalizaciones y espectrogramas}

Una vez generado el dataset, este debe ser preprocesado para que la red lo pueda manejar correctamente.
Así pues, se extrae la magnitud de la señal, se convierte a decibelios (dB) y se guarda como tensores (formato \textit{.pt}). Además,
se genera un archivo (spec\_mode.txt) que indica el tipo de espectrograma utilizado.

Para terminar con el tratamiento del dataset, se estructuran los tensores en una clase Dataset de PyTorch
(\textit{SpectrogramTensorDataset}). Se lee el registro de etiquetas (\textit{labels.cvs}), el cual contiene todos los
valores de los parámetros de cada audio, y se calcula la media y desviación estándar global de cada uno de ellos. Con estas estadísticas
se aplica una normalización de puntuación estándar (\textit{Z-score}), nivelando así las diferencias entre las escalas
de los parámetros, consiguiendo igualar las importancias de cada uno en el cálculo del gradiente.
Finalmente, el dataset se divide, mediante una semilla fija, en: 70\% para entrenamiento, 15\% validación y se guarda otro 15\%
como conjunto de pruebas (\textit{Test set}).

\section{Diseño de la Arquitectura CNN}
\label{sec:entrenamiento_final}
Con el fin de experimentar con distintos modelos en el entrenamiento, se proponen dos arquitecturas distintas.

\begin{itemize}
    \item \textbf{CNNRegressorSimple: }Está formada por un \textit{Encoder} de tres bloques convolucionales,
    un \textit{Bottleneck} y una cabeza de regresión pura (\textit{Global Average Pooling}). 

    Primero se pasa por un Codificador (\textbf{Encoder}), el cual consta de una secuencia de tres bloques
    convolucionales que convierten el espectrograma inicial en una representación comprimida, extrayendo así sus 
    características fundamentales. Para ello, cada bloque aplica núcleos (kernels) de 3x3. Se pasa por una capa de normalización por lotes (Batch Normalization), que estabiliza y acelera el entrenamiento, y 
    una activación no lineal ReLU. Finalmente, se aplica una reducción espacial en cada bloque (Max Pooling).
    Con todo esto, la matriz se ha reducido a una octava parte de su tamaño original y la profundidad de la red
    ha incrementado de 1 a 128 canales. Como resultado, se obtiene una representación comprimida fiel al sonido original.

    Al salir del \textbf{Encoder}, el flujo de datos atraviesa lo que se denomina como
    cuello de botella (\textbf{Bottleneck}). Este bloque convolucional no
    aplica una reducción espacial, manteniendo estables las dimensiones de la matriz. Su función es duplicar el número
    de canales (profundidad del tensor), pasando de 128 a 256 canales. Esto se hace con la finalidad de 
    mezclar los patrones que el \textbf{Encoder} ha extraído previamente, creando una representación
    densa y unificada.

    Después se pasa por la fase de Inferencia Numérica, donde ocurre la regresión de los parámetros del sintetizador FM. Se aplica una capa de agrupamiento global (Adaptive Average Pooling) ya que los 
    parámetros describen el sonido completo y no un punto temporal concreto en el espectrograma.
    Para ello, se calcula la media de todas las posiciones de cada canal, quedándose con un único
    número como resumen por canal. Finalmente, el vector unidimensional, resultado del proceso anterior,
    atraviesa un \textbf{Perceptrón Multicapa}, que lo proyecta hacia las 8 neuronas de salida.
    
    A pesar de que la red ejecuta un entrenamiento rápido, al carecer de reconstrucción 
    no se garantiza que el \textit{Encoder} esté extrayendo de manera acertada la estructura del espectrograma, prediciendo únicamente números.
        
    \item \textbf{CNNRegressor5: }Se trata de una expansión del modelo anterior, incorporando una segunda cabeza de predicción. Se añade la fase de Reconstrucción (\textbf{Decoder}), la cual se encarga de reconstruir 
    el espectrograma original a partir de la representación comprimida del \textbf{Bottleneck}. 
    Como el \textbf{Encoder} comprime tres veces, el \textbf{Decoder} descomprime tres veces. Para ello,
    utiliza tres capas convolucionales transpuestas (\textit{ConvTranspose2d}).

    El motivo para utilizar el \textbf{Decoder} es como método de regularización forzada del \textbf{Encoder}.
    En el caso en el que el codificador borre información importante, como armónicos o patrones temporales, que no afecte directamente a los parámetros
    de la regresión, el \textbf{Decoder} no será capaz de redibujar correctamente el espectrograma, disparando así
    el error en la función de pérdida híbrida. Aparte, si el \textbf{Bottleneck} no guarda suficiente información, ocurrirá lo mismo.
    Todo esto fuerza, mediante el cálculo de gradientes, a que las capas iniciales aprendan a conservar el timbre y la representación
    del sonido.
\end{itemize}

\section{Funciones de Pérdida} %falta hablar de los papers

\subsection{El problema de la inyectividad}
La busqueda de una función de pérdida se vio totalmente condicionada por una característica esencial de la síntesis FM: su inyectividad es nula. Esto quiere decir que diferentes combinaciones de parámetros pueden dar lugar a sonidos idénticos.

Una red implementada con una función de pérdida que trate de predecir los parámetros exactos del audio original, puede verse afectada negativamente por esta característica, cosa que se comprobó en un primer prototipo cuya pérdida se media de esta forma utilizando el Error Cuadrático Medio (\textit{MSELoss}). Un ejemplo que sufría este modelo, radical pero que explica muy bien este resultado, es un caso en el que el índice de modulación sea 0. En ese caso no es verdaderamente relevante el valor del ratio, sin embargo, 
la función de pérdida, implementada con \textit{MSE}, penaliza gravemente que este difiera del original. 

Otro ejemplo, más común, ocurre en los casos en los que se trabaja sobre un índice de modulación alto, de tal forma que la frecuencia moduladora predomina sobre la portadora. Si tenemos un \textit{Carrier} de 1000 Hz y un \textit{Ratio} de 1.0, la frecuencia moduladora resultante será de 1000 Hz, generándose los armónicos en saltos de mil en mil hercios. Por otro lado, con un \textit{Carrier} de 2000 Hz y un \textit{Ratio} de 0.5, se obtiene la misma moduladora de 1000 Hz y, por lo tanto, los mismos armónicos. Todo esto sumado a un índice de modulación alto, harían que estos dos audios, paramétricamente muy distantes, suenen muy similares al oído humano. 

\subsection{Funciones}

Manteniendo el objetivo de ampliar el espacio de pruebas, se permite también el uso de diferentes funciones de pérdida.

El modelo simple, evalua la diferencia paramétrica únicamente mediante
\textbf{SmoothL1 Loss}. Esta se utiliza por su robustez y versatilidad en la regresión numérica: cuando la diferencia 
entre el valor real y el predicho es muy pequeña, se comporta como un Error Cuadrático (L2), lo que suaviza la convergencia
de gradientes; y ante desviaciones grandes, se asemeja al Error Absoluto (L1), de forma que los casos aislados que se alejan del resto 
no desestabilicen la red. Esta función de pérdida tiene la desventaja de no tener en cuenta la no-inyectividad de la síntesis FM.

En cuanto al modelo CNNRegressor5, se permite la elección de una de las siguientes funciones de pérdida:
\begin{itemize}
    \item \textbf{HybridLoss: }
    Es una función de pérdida híbrida, que combina el Error Absoluto Medio (L1) de la reconstrucción del 
    espectrograma, un refuerzo estructural de Convegencia Espectral (SC), y la penalización paramétrica anterior.
    La primera se lleva el peso principal (de 1.0), calcula el Error Absoluto Medio píxel a píxel entre el espectrograma original y el reconstruido (en decibelios).
    La Convergencia Espectral (SC) se aplica con un peso medio (de 0.5) para evitar los espectros borrosos generados por L1. Esta métrica calcula la norma de Frobenius 
    para evaluar la energía global (relación armónica), regularizando el timbre. A la penalización paramétrica se le da una importancia muy baja (0.05), ya que no es demasiado importante
    debido al carácter no inyectivo de la síntesis FM.
        
    \item \textbf{MultiScaleSpectralLoss: } Esta función está inspirada en el modelo de síntesis diferencial DDSP. Al estar utilizando 
    espectrogramas, se simulan los distintos tamaños de ventana de Fourier mediante agrupamientos espaciales temporales (\textit{Average Pooling}).
    Se agrupa en seis factores (de 1, 2, 4, 8, 16 y 32) y se mide el error en cada una de esas escalas, siendo la media aritmética de esos errores
    el valor de pérdida total. El cálculo del error de cada escala primero se calcula la diferencia sobre las magnitudes
    lineales, la cual es útil para capturar picos de energía que pueden ser breves. A esto se le suma después la diferencia sobre las magnitudes en decibelios (logarítmico),
    multiplicada por un factor $\alpha$ de 1.0, por lo que se le está danto la misma relevancia que el error lineal, capturando así perfectamente
    la envolvente global y los momentos de baja energía. Mediante esta función, se consigue compara el sonido obtenido del original
    en varios niveles de resolución simultáneamente. Por último, se le suma el error paramétrico (\textit{SmoothL1}) con un peso muy bajo (0.05) con el objetivo
    de que la red no deje de tener mínimamente en cuenta los valores de los parámetros.   
\end{itemize}


\subsection{Modelo de Pruebas simplificado}

Se decide mantener un modelo simplificado, para la realización de pruebas, que infiere únicamente 3 parámetros: frecuencia portadora (\textit{carrier}), frecuencia moduladora (\textit{ratio}) e índice de modulación (\textit{index}). 

La topología de la red (\textit{CNNRegressor4}) es igual a la compleja del modelo principal (\textit{CNNRegressor5}), solo que la salida son 3 parámetros en vez de 8. La función de pérdida es la \textit{HybridLoss} explicada en el apartado anterior. 

